\documentclass[lettersize,journal]{IEEEtran}
\usepackage{amsmath,amsfonts}
\usepackage{algorithmic}
\usepackage{algorithm}
\usepackage{array}
\usepackage[caption=false,font={footnotesize,rm}, labelformat=simple]{subfig}
% 1. Add 'labelformat=simple' to stop the caption from adding EXTRA parentheses

% 2. Manually add parentheses to the subfigure counter itself.
%    This forces the \ref command to include them (e.g., "1(a)"),
%    and because we used 'simple' above, the caption stays as "(a)".
\renewcommand{\thesubfigure}{(\alph{subfigure})}
\usepackage{textcomp}
\usepackage{stfloats}
\usepackage{url}
\usepackage{verbatim}
\usepackage{graphicx}
\usepackage{cite}
\usepackage{cleveref}
\usepackage{multirow}

\usepackage{pgfplots}
\pgfplotsset{compat=1.18}
\usepgfplotslibrary{groupplots} % Essential for the "groupplot" environment
\raggedbottom
\hyphenation{op-tical net-works semi-conduc-tor IEEE-Xplore}

% Updated with editorial comments based on text dump provided

\begin{document}

\title{Design and Implementation of a Non-coupled Integrated Transformer for a Two-phase Interleaved Flyback PFC Rectifier for USB PD 3.1 Applications}

\author{Karl Raymond Dela Cruz Roque, \IEEEmembership{Student Member, IEEE}, Haruki Hirasawa, Sihoon Choi, \IEEEmembership{Member, IEEE}, Yu Yonezawa, \IEEEmembership{Member, IEEE}, Jun Imaoka, \IEEEmembership{Member, IEEE}, and Masayoshi Yamamoto, \IEEEmembership{Member, IEEE}
        % <-this % stops a space
\thanks{Manuscript received Month xx, 2xxx; revised Month xx, xxxx; accepted Month x, xxxx. (Corresponding author: Sihoon Choi).}
\thanks{K. R. D. Roque, S. Choi, Y. Yonezawa, J. Imaoka and M. Yamamoto are with the Institute of Materials and Systems for Sustainability, Nagoya University, Nagoya, 464-8601, Japan (e-mail: roque.karl.raymond.dela.cruz.g46@f.mail.nagoya-u.ac.jp; choi.sihoon.w4@f.mail.nagoya-u.ac.jp; yonezawa.yu.c9@f.mail.nagoya-u.ac.jp; imaoka.jun.n9@f.mail.nagoya-u.ac.jp; yamamoto.masayoshi.n5@f.mail.nagoya-u.ac.jp).}
\thanks{H. Hirasawa is with [Affiliation Needed], Japan (e-mail: ).}}

% The paper headers
\markboth{Journal of Power Electronics,~Vol.~XX, No.~X, Month~202X}%
{Roque \MakeLowercase{\textit{et al.}}: Design and Implementation of a Non-coupled Integrated Transformer}

\maketitle

\begin{abstract}
Universal serial bus power delivery (USB PD 3.1) is a widely-applied standard for electronic device charging. In 2021, USB PD 3.1 was announced with a variable output of 5 V to 48 V and a maximum output power of 240 W. A multi-phase flyback rectifier is capable of power factor correction, galvanic isolation, and wide output voltage range, and is proposed as an alternative topology for USB PD. However, the number of magnetic components increases with the number of phases, which can lead to lower power density. This paper proposes the design of a non-coupled integrated transformer for a two-phase interleaved flyback rectifier to reduce the size and number of magnetic components. The independent operation of the integrated transformer in two phases was verified using a 240-W two-phase interleaved flyback PFC rectifier at 100 kHz switching frequency operating at discontinuous conduction mode (DCM), both in simulation and experiment. After conducting a volume analysis through the area product method, the proposed transformer is up to 56\% smaller compared to its conventional counterpart.
\end{abstract}

\begin{IEEEkeywords}
Integrated magnetics, Transformer design, Two-phase interleaved flyback rectifier, USB PD 3.1, PFC, Core volume analysis.
\end{IEEEkeywords}

\section{Introduction}
\IEEEPARstart{U}{niversal} serial bus power delivery (USB PD) is a technology used for convenient, portable, fast charging in electronic devices. The latest standard, USB PD 3.1, was announced in 2021, with the maximum power supplied being up to 240 W, allowing for higher power applications \cite{ref1}. Adapters for USB PD 3.1 must be capable of power factor correction (PFC) for supporting 75 W and above, be equipped with galvanic isolation for user safety, and be capable of providing an output voltage of 5 V to 48 V \cite{ref2}.

A flyback rectifier has been widely used for commercial adapters due to its simple configuration, easy PFC control, and wide output voltage range. Improved configurations of flyback rectifiers such as active clamp flyback converters \cite{ref6, ref7} and bridgeless flyback rectifiers have also been proposed for higher converter efficiency \cite{ref8}, \cite{ref9}, \cite{ref10}, \cite{ref11}. However, they are generally more practical for lower power applications due to the inherent high current stress on the switches by these topologies. 

Several power supply configurations for USB PD 3.1 have been proposed to meet the requirements. One example is a combination of a PFC rectifier, an LLC converter, and a buck converter as shown in Fig. \ref{fig:1a} \cite{ref3, ref4}. While this configuration achieves high power density, the additional buck stage required to compensate for the limited gain range of the LLC increases the total component count and system cost. An asymmetrical half bridge (AHB) flyback converter in the second stage of a boost PFC circuit was introduced as shown in \ref{fig:1b} \cite{ref5}. The AHB allows for zero-voltage switching (ZVS) by the application of resonant capacitors in series placed on the primary side of the transformer. As a single-phase topology, the AHB processes the entire load power through a single transformer and switch pair. While this minimizes the component count, it inherently concentrates thermal stress in specific circuit elements.

\begin{figure}[htbp]
\centering
% (a) Top circuit
        \subfloat[]{
        \includegraphics[width=\columnwidth]{Figures_v2/Fig1a.pdf} % Make sure this file exists
        \label{fig:1a}
        }
        \\ % This double backslash forces the next image to the line below (vertical stack)
        % (b) Middle circuit
        \subfloat[]{
        \includegraphics[width=\columnwidth]{Figures_v2/Fig1b.pdf}
        \label{fig:1b}
        }
        \\ % Line break
        % (c) Bottom circuit
        \subfloat[]{
        \includegraphics[width=\columnwidth]{Figures_v2/Fig1c.pdf}
        \label{fig:1c}
        }
        \caption{Power supply configurations proposed for PD 3.1 applications: (a) PFC boost, LLC, and buck converter combination in \cite{ref3,ref4}, (b) PFC boost with AHB flyback \cite{ref1}, (c) two-phase interleaved flyback PFC rectifier.}
        \label{fig:configurations}
\end{figure}

Interleaving rectifiers have been considered to increase the maximum output power. Two-phase interleaving allows half of the total power to be supplied in each phase, therefore reducing the possible stress on the circuit elements \cite{ref12}. The topology is shown in Fig. \ref{fig:1c}, where the two phases are parallel in the input and output side. The source nodes of switches $S_1$ and $S_2$ are common, and only one power supply is required for gate drive circuits. A common cathode package can be used for diodes $D_1$ and $D_2$, therefore fewer components are used for the circuit. 

Operating multi-phase flyback converters across a wide output range presents some challenges. For instance, to ensure high efficiency across wide-range applications, converters must commonly employ phase-shedding, for which various control schemes have been proposed \cite{Zhang_2013}, \cite{Su_2013, Chen_2014}. This strategy strategically deactivates specific operational phases depending on the load condition to optimize light-load efficiency. However, the implementation of these control schemes often necessitates complicated designs and additional active components, thereby increasing the complexity of the circuit operation. 

Moreover, increasing the number of phases inherently means an increase in the magnetic components required for the circuit. In the topology utilized for this study, although fewer active components are used overall, two discrete transformers are required, which compromises the power density of the converter. Integrated magnetics have been proposed and widely studied to improve the power density of various types of converters. Because integrated magnetics consolidate multiple discrete components through separate windings on a single magnetic core, these components inherently share the same flux paths. Consequently, the winding structure must be carefully designed to control the magnetic flux. Depending on the winding structure and the utilized flux paths, once mutual coupling is present between windings, the phases inextricably influence each other. In single-phase operation, when the phases are highly coupled, the active phase inevitably induces a voltage across the inactive phase, driving massive unwanted circulating currents that severely degrade light-load efficiency \cite{Yang_2020, Z_Li_2025}. Even during heavy-load conditions where multiple phases are in operation, coupling can still enable unintended power delivery across the non-operational phase intervals, leading to input current distortion and additional induced losses, thereby reducing overall efficiency \cite{Yang_2020}.

Conversely, for strictly heavy-load applications, the design priority changes. Strong magnetic coupling passively forces current equalization and minimizes capacitive ripple, while the integrated magnetics further contribute to improving power density. This intentional heavy-load coupling is highly advantageous and prominently utilized by Ahmad et al., Dou et al., and Nakakohara et al. for fixed-load server power supplies and LLC resonant converters \cite{Ahmad}, \cite{Dou, Nakakohara}. Furthermore, modern interleaved flybacks still require extra active switching devices for clamp circuits to guarantee this auto-sharing \cite{Z_Chen_s42835}. Ultimately, while strongly coupled magnetic structures are highly effective for fixed heavy-load design targets, they are fundamentally incompatible with the strict independent phase operation required by wide-range, variable-power converters. Lastly, other integrated magnetic structures exist for alternative targets, such as noise reduction through active cancellation within the windings.

Therefore, to make integrated magnetics more compatible with the variable-load requirements of USB-PD 3.1, approaches to decoupling phases have been explored in terms of the magnetic structure to bypass these circulating currents, aside from relying solely on control-based phase shedding. An approach proposed a Four Quadrants Integrated Transformer (FQIT) for a two-phase converter to achieve perfectly decoupled phases by physically separating the flux paths using a complex 3-D orthogonal-flux core geometry \cite{Ouyang}. However, such a structure utilizes a custom, non-standard core, which constrains commercial manufacturability and increases fabrication costs. Similarly, an integrated nonlinear coupled inductor designed to naturally decouple four phases at light loads has been proposed \cite{Z_Li_2025}. However, achieving this variable coupling relies heavily on highly customized, multi-piece L- and T-shaped core assemblies and geometrically self-adjusting structures, which introduce tight assembly tolerances and significant manufacturing complexity \cite{Jin_2024}.

Therefore, to address these limitations, this paper proposes a non-coupled integrated transformer for a two-phase interleaved flyback rectifier as shown in Fig. \ref{fig:2a}. To ensure independent phase operation without magnetic interference, the structure is designed to achieve zero magnetic coupling between phases. The proposed winding structure is shown in Fig. \ref{fig:2b} where $N_{12}$ is the number of turns between terminals 1 and 2. Half of $N_{12}$ and $N_{34}$ are wound at each outer leg with opposite polarity, so as to allow for flux cancellation at the center leg. On the other hand, $N_{56}$ and $N_{78}$ are also wound with opposite polarity on the center leg, and allows for the transformer $N_{12}$ and $N_{34}$ to not operate during its operation. 

\begin{figure}[b]
\centering
% (a) Top circuit
        \subfloat[]{
        \includegraphics[width=\columnwidth]{Figures_v2/Fig2a.pdf} % Make sure this file exists
        \label{fig:2a}
        }
        \\ % This double backslash forces the next image to the line below (vertical stack)
        % (b) Middle circuit
        \subfloat[]{
        \includegraphics[width=0.6\columnwidth]{Figures_v2/Fig2b.pdf}
        \label{fig:2b}
        }
        \caption{The two-phase interleaved flyback rectifier (a), with the transformer and its winding structure in (b).}
\end{figure}

The zero-coupling coefficient between the two separate transformers simplifies implementation of rectifiers and transformer design. The operation of the rectifiers in each interleaved phase still follows that of a conventional flyback converter, preserving simplicity in control and timing. The design is also highly dependent on the duty cycle, an easily adjustable parameter in circuit design. The converter in this study aims to operate in DCM in order to allow for easier implementation of PFC and to conform to the standards of USB PD 3.1. It also realizes independent phase operation, allowing the circuit to function in one-phase mode if desired, thereby enhancing flexibility in circuit operation and allowing simple implementation of phase-shedding. Furthermore, this entirely decoupled winding structure is easily achievable using a conventional 2-D EE-core, successfully bypassing the manufacturing complexities of custom geometries.

The rest of this paper is organized as follows: Section II presents an operational analysis of a two-phase interleaved flyback rectifier with the integrated transformer, along with its design method. Section III describes the size reduction effect of the transformer. Sections IV and V discuss the results of the simulation and experimental verification respectively, and Section VI presents the conclusions of this paper.

\section{Design Method of the Proposed Transformer}

\subsection{Operation Analysis of the Proposed Flyback Rectifier and Transformer}
The proposed integrated transformer, shown in Fig. \ref{fig:2b} was designed for a two-phase interleaved flyback PFC rectifier, shown in Fig. \ref{fig:2a} operating in discontinuous conduction mode (DCM). The switches ($S_1$, $S_2$ in Fig. \ref{fig:2a}) operate with a fixed duty cycle $D$ and a 180-degree phase shift.

Half of $N_{12}$ and $N_{34}$ are wound at each outer leg with opposite polarity. Moreover, $N_{56}$ and $N_{78}$ are also wound with opposite polarity at the center leg. The transformer utilized an EE Core with a center leg that has a cross sectional area twice its outer leg, as per commercial standard. The magnetizing inductances for both transformers $N_{12}-N_{34}$ and $N_{56}-N_{78}$ were set to be equal, and are denoted as $L_{\mathrm{m}_1}$ and $L_{\mathrm{m}_2}$ respectively. The saturation magnetic flux in each separate transformer during each phase is expected to be consistent.

The operational principle of the rectifier with the proposed transformer is illustrated in Fig. \ref{fig:operation_modes}. Additionally, Fig. \ref{fig:4} presents the key waveforms of the six distinct operation modes. The rectifier is designed with duty cycle $D \le 0.5 \frac{n V_{\mathrm{o}}}{V_{\mathrm{in}_{\mathrm{max}}}}$ with $n$ as the turns ratio, $V_{\mathrm{o}}$ as the output voltage, and $V_{\mathrm{in}_{\mathrm{max}}}$ as the peak value of the input voltage. This condition ensures the magnetizing current of one phase returns to zero before the other phase peaks. This duration is defined as the discharge time $t_{\mathrm{dis}}$ and is given in terms of circuit parameters by \eqref{eq:discharge}:

\begin{figure*}[htbp]
    \centering
    
    % --- LEFT COLUMN (Modes 1, 2, 3) ---
    % These will automatically get labels (a), (b), (c)
    \begin{minipage}{0.48\textwidth}
        \centering
        
        \subfloat[\label{fig:mode_1}]{%
            \includegraphics[width=\linewidth]{Figures_v2/Fig3a.pdf}%
        }
        
        \vspace{10pt} % Spacing between Mode 1 and 2
        
        \subfloat[\label{fig:mode_2}]{%
            \includegraphics[width=\linewidth]{Figures_v2/Fig3c.pdf}%
        }
        
        \vspace{10pt} % Spacing between Mode 2 and 3
        
        \subfloat[\label{fig:mode_3}]{%
            \includegraphics[width=\linewidth]{Figures_v2/Fig3e.pdf}%
        }
    \end{minipage}
    \hfill % Horizontal space between the two columns
    % --- RIGHT COLUMN (Modes 4, 5, 6) ---
    % These will automatically get labels (d), (e), (f)
    \begin{minipage}{0.48\textwidth}
        \centering
        
        \subfloat[\label{fig:mode_4}]{%
            \includegraphics[width=\linewidth]{Figures_v2/Fig3b.pdf}%
        }
        
        \vspace{10pt} 
        
        \subfloat[\label{fig:mode_5}]{%
            \includegraphics[width=\linewidth]{Figures_v2/Fig3d.pdf}%
        }
        
        \vspace{10pt} 
        
        \subfloat[\label{fig:mode_6}]{%
            \includegraphics[width=\linewidth]{Figures_v2/Fig3f.pdf}%
        }
    \end{minipage}
    
    \caption{Operational principle of the two-phased interleaved flyback rectifier with the proposed integrated transformer at $D \le 0.5 \frac{n V_{\mathrm{o}}}{V_{\mathrm{in}_{\mathrm{max}}}}$. There are six distinct operational modes shown: (a) Mode 1 ($t_1 < t < t_2$), (b) Mode 2 ($t_2 < t < t_3$), (c) Mode 3 ($t_3 < t < t_4$), (d) Mode 4 ($t_4 < t < t_5$), (e) Mode 5 ($t_5 < t < t_6$), and (f) Mode 6 ($t_6 < t < t_7$).}
    \label{fig:operation_modes}
\end{figure*}

\begin{equation}
\label{eq:discharge}
t_{\mathrm{dis}} = \frac{V_{\mathrm{in}_{\mathrm{max}}} DT_{\mathrm{s}}}{n V_{\mathrm{o}}}
\end{equation}
with $T_{\mathrm{s}}$ as the switching period. For the condition given by $D$ in this analysis, $t_{\mathrm{dis}}$ corresponds to the interval between $t_3$ and $t_5$ in Fig. \ref{fig:4}.

\begin{figure}[htbp]
\centering
\includegraphics[width=\columnwidth]{Figures_v2/Fig4.pdf}

\caption{Key waveforms in the operation of the two-phase interleaved flyback rectifier with the proposed transformer, with $D \le 0.5 \frac{n V_{\mathrm{o}}}{V_{\mathrm{in}_{\mathrm{max}}}}$. }
\label{fig:4}
\end{figure}

To describe the magnetic state of the core, the superposition principle is used, as shown in Fig. \ref{fig:5a}. The magnetic flux density component $B'$ is generated by the positive cycle transformer defined as $N_{12}-N_{34}$. The magnetic flux density component $B''$ represents the component generated by the negative cycle transformer defined as $N_{56}-N_{78}$. The total flux in any leg is the algebraic sum of these components. This superposition determines the peak flux density required for the core design, which is analyzed in detail in Section II-B.

Modes 1, 2, and 3 in Figs. \ref{fig:mode_1}-(c) cover the operation during the positive AC input into the rectifier, while Modes 4, 5, and 6 in Figs. \ref{fig:mode_4}-(f) describe the operation during the negative AC input. For the following modes, $B'_{\mathrm{o}_1}$, $B'_{\mathrm{o}_2}$ and $B'_{\mathrm{c}}$ represent the magnetic flux densities in the respective left, right, and center legs of the core caused by the transformer active in positive AC input $N_{12}-N_{34}$, respectively, which are shown visually in Fig. \ref{fig:5b}. On the other hand, $B''_{\mathrm{o}_1}$, $B''_{\mathrm{o}_2}$ and $B''_{\mathrm{c}}$ represent the magnetic flux densities in the respective left, right, and center legs of the core caused by the transformer active in negative AC input $N_{56}-N_{78}$, respectively, which are detailed visually in Fig. \ref{fig:5c}.

The detailed analysis of each operating mode is as follows:


\begin{enumerate}
    \item Mode 1 ($t_1 < t < t_2$, Fig. \ref{fig:mode_1}): At time $t_1$, switch $S_1$ turns on and $S_2$ remains off. The transformer $N_{12}-N_{34}$ begins storing energy in its magnetizing inductance. Its magnetizing current $i_{L_{\mathrm{m}_1}}$ and the corresponding flux densities in the outer legs $B'_{\mathrm{o}_1}$ and $B'_{\mathrm{o}_2}$ linearly increase from zero. Simultaneously, transformer $N_{56}-N_{78}$ completes its discharge phase, with its magnetizing current $i_{L_{\mathrm{m}_2}}$ and secondary diode current $i_{\mathrm{D}_2}$ decreasing to zero at $t_2$.

    \item Mode 2 ($t_2 < t < t_3$, Fig. \ref{fig:mode_2}): Mode 2 starts right after transformer $N_{56}-N_{78}$ completes its discharging phase, which means its associated currents $i_{\mathrm{D}_2}$, $i_{L_{\mathrm{m}_2}}$, and flux densities $B''_{\mathrm{o}_1}$, $B''_{\mathrm{o}_2}$, and $B''_{\mathrm{c}}$ are zero. Now, with $S_1$ still on, transformer $N_{12}-N_{34}$ continues storing energy until $i_{L_{\mathrm{m}_1}}$ reaches its peak value at $t_3$. The peak primary current in the positive cycle $i_{\mathrm{S}_{1_{\mathrm{pk}}}}$ is given by \eqref{eq:1}:
    \begin{equation}
    \label{eq:1}
    i_{\mathrm{S}_{1_{\mathrm{pk}}}} = \frac{V_{\mathrm{in}} \cdot DT_{\mathrm{s}}}{L_{\mathrm{m}_1}}
    \end{equation}
    where $V_{\mathrm{in}}$ is the input voltage, $D$ is the duty cycle, $T_{\mathrm{s}}$ is the switching period, and $L_{\mathrm{m}_1}$ is the magnetizing inductance of transformer $N_{12}-N_{34}$. Consequently, the flux densities $B'_{\mathrm{o}_1}$ and $B'_{\mathrm{o}_2}$ also reach their peak values $B'_{\mathrm{o}_{\mathrm{pk}}}$, which are given by \eqref{eq:2}:
    \begin{equation}
    \label{eq:2}
    B'_{\mathrm{o}_{\mathrm{pk}}} = \frac{L_{\mathrm{m}_1} i_{\mathrm{S}_{1_{\mathrm{pk}}}}}{N_{12} A_{\mathrm{o}}}
    \end{equation}
    with $N_{12}$ as the number of turns in the primary side of the active transformer and $A_{\mathrm{o}}$ is the cross-sectional area of the outer leg, where the transformer $N_{12}-N_{34}$ is wound.

    \item Mode 3 ($t_3 < t < t_4$, Fig. \ref{fig:mode_3}): Mode 3 starts right after transformer $N_{12}-N_{34}$ finishes charging. At $t_3$, Switch $S_1$ turns off, initiating the discharge phase of transformer $N_{12}-N_{34}$. The magnetizing current $i_{L_{\mathrm{m}_1}}$, diode current $i_{\mathrm{D}_1}$, and flux densities $B'_{\mathrm{o}_1}$ and $B'_{\mathrm{o}_2}$ begin to decrease linearly. Transformer $N_{56}-N_{78}$ remains inactive in this period.
\end{enumerate}

In Modes 1-3, the magnetomotive force and magnetic resistance of both outer legs are equal. The magnetic flux in the center leg is canceled. This causes no electromotive force generated in transformer $N_{56}-N_{78}$. Therefore, the transformer formed by the outer legs operates independently from the transformer in the center leg. The integrated transformer is equivalent to that of Fig. \ref{fig:5b} and its equivalent magnetic circuit is shown in Fig. \ref{fig:6}a.

\begin{enumerate}
    \setcounter{enumi}{3} % Ensures the numbering continues from 4
    \item Mode 4 ($t_4 < t < t_5$, Fig. \ref{fig:mode_4}): Mode 4 operates similarly to the positive cycle. Switch $S_2$ turns on at $t_4$ and $S_1$ remains off. Transformer $N_{56}-N_{78}$ stores energy in its magnetizing inductance. Its current $i_{L_{\mathrm{m}_2}}$ and the corresponding flux densities $B''_{\mathrm{o}_1}$, $B''_{\mathrm{c}}$, and $B''_{\mathrm{o}_2}$ linearly increase. Simultaneously, transformer $N_{12}-N_{34}$ completes its discharge phase, with $i_{L_{\mathrm{m}_1}}$ and $i_{\mathrm{D}_1}$ decreasing to zero at $t_5$.

    \item Mode 5 ($t_5 < t < t_6$, Fig. \ref{fig:mode_5}): Mode 5 starts after the discharge of transformer $N_{12}-N_{34}$ completes. Now, with Switch $S_2$ still on, transformer $N_{56}-N_{78}$ continues storing energy until $i_{L_{\mathrm{m}_2}}$ reaches its peak value at $t_6$. The peak primary current in the negative cycle $i_{\mathrm{S}_{2_{\mathrm{pk}}}}$ is given by \eqref{eq:3}:
    \begin{equation}
    \label{eq:3}
    i_{\mathrm{S}_{2_{\mathrm{pk}}}} = \frac{V_{\mathrm{in}} \cdot DT_{\mathrm{s}}}{L_{\mathrm{m}_2}}
    \end{equation}
    where $L_{\mathrm{m}_2}$ is the magnetizing inductance of transformer $N_{56}-N_{78}$. Consequently, the flux densities $B''_{\mathrm{o}_1}$, $B''_{\mathrm{c}}$, and $B''_{\mathrm{o}_2}$ also reach their peak values. The peak value of $B''_{\mathrm{c}_{\mathrm{pk}}}$ is given by \eqref{eq:4}:
    \begin{equation}
    \label{eq:4}
    B''_{\mathrm{c}_{\mathrm{pk}}} = B''_{\mathrm{o}_{\mathrm{pk}}} = \frac{L_{\mathrm{m}_2} i_{\mathrm{S}_{2_{\mathrm{pk}}}}}{N_{56} A_{\mathrm{c}}}
    \end{equation}
    with $N_{56}$ as the number of turns in the primary side of the active transformer and $A_{\mathrm{c}}$ is the cross-sectional area of the center leg where the transformer $N_{56}-N_{78}$ is wound.

    \item Mode 6 ($t_6 < t < t_7$, Fig. \ref{fig:mode_6}): Switch $S_2$ turns off, initiating the discharge phase of transformer $N_{56}-N_{78}$. The currents $i_{L_{\mathrm{m}_2}}, i_{\mathrm{D}_2},$ and flux densities $B''_{\mathrm{o}_1}$, $B''_{\mathrm{c}}$, and $B''_{\mathrm{o}_2}$ begin to decrease linearly. Transformer $N_{12}-N_{34}$ remains inactive.
\end{enumerate}

For Modes 4-6, since the magnetic resistance of both outer legs is equal, the magnetic flux is divided equally into both outer legs. Furthermore, since the number of turns of both outer legs is equal and wound in an opposite polarity, the electromotive force generated in $N_{12}-N_{34}$ is canceled. Therefore, the transformer formed by the center leg operates independently from the transformer in the outer leg, and the integrated transformer is equivalent to that of Fig. \ref{fig:5c}. The equivalent magnetic circuit for this case is shown in Fig. \ref{fig:6}b.

The proposed winding structure achieves zero magnetic coupling between the two transformers. Consequently, the distinct magnetic circuits share a single core without magnetic interference. Each transformer maintains independent operation despite the physical integration. This structure eliminates the need for two discrete magnetic components. As a result, the total volume of the magnetic components is significantly reduced.

\begin{figure}[b!]
\centering
% (a) Top circuit
        \subfloat[]{
        \includegraphics[width=0.6\columnwidth]{Figures_v2/Fig5a.pdf} % Make sure this file exists
        \label{fig:5a}
        }
        \\ % This double backslash forces the next image to the line below (vertical stack)
        % (b) Middle circuit
        \subfloat[]{
        \includegraphics[width=0.6\columnwidth]{Figures_v2/Fig5b.pdf}
        \label{fig:5b}
        }
        \\
        \subfloat[]{
        \includegraphics[width=0.53\columnwidth]{Figures_v2/Fig5c.pdf}
        \label{fig:5c}
        }
        \caption{The proposed integrated transformer during the two phases of the converter, with the flux density components in each leg illustrated in (a). (b) is realized during the positive AC cycle, while (c) is realized during the negative cycle.}
\end{figure}
\begin{figure}[b!]
\centering
% (a) Top circuit
        \includegraphics[width=0.87\columnwidth]{Figures_v2/Fig6.pdf} % Make sure this file exists
        \label{fig:6}
        \caption{Equivalent magnetic circuits for the proposed transformer. (a) shows the equivalent magnetic circuit for $N_{12}-N_{34}$, while (b) shows the equivalent magnetic circuit for $N_{56}-N_{78}$.}
\end{figure}

\subsection{Saturation Magnetic Flux Density Analysis}
\begin{figure*}[htbp]
    \centering
    % The figure spans two columns, so we use figure*
    
    % --- (a) Case 1 ---
    % \subfloat[] automatically generates the "(a)" label. 
    % We leave the bracket empty [] because the text "Case 1" is in the main caption, not the sub-label.
    \subfloat[\label{fig:case1}]{%
        \includegraphics[width=0.32\textwidth]{Figures_v2/Fig7a.pdf}%
    }
    \hfill % Adds flexible space between images
    % --- (b) Case 2 ---
    \subfloat[\label{fig:case2}]{%
        \includegraphics[width=0.32\textwidth]{Figures_v2/Fig7b.pdf}%
    }
    \hfill
    % --- (c) Case 3 ---
    \subfloat[\label{fig:case3}]{%
        \includegraphics[width=0.32\textwidth]{Figures_v2/Fig7c.pdf}%
    }
    
    \caption{Theoretical, superimposed magnetic flux density waveforms in the integrated transformer during circuit operations: (a) Case 1, (b) Case 2, and (c) Case 3.}
    \label{fig:flux_waveforms}
\end{figure*}

The circuit operation is classified into three patterns determined by $D$ and the boundary condition $D \le 0.5 \frac{n V_{\mathrm{o}}}{V_{\mathrm{in}_{\mathrm{max}}}}$. This classification, illustrated in Fig. 7, arises because the integrated single-core structure permits simultaneous magnetizing current flow in multiple active windings, resulting in distinct flux superposition compared to conventional discrete transformers. The manner of superposition due to the magnetizing current can be controlled by $t_{\mathrm{dis}}$. Consequently, the saturation magnetic flux density $B_{\mathrm{max}}$ is a critical parameter for the transformer geometry. It can be expressed in terms of the saturation magnetic flux density in each phase, denoted as $B'_{\mathrm{max}}$ and $B''_{\mathrm{max}}$, with $B'_{\mathrm{max}} = B''_{\mathrm{max}}$ since $L_{\mathrm{m}_1} = L_{\mathrm{m}_2}$. The flux density in each leg is examined further using the principle of superposition, with the expressions of the superimposed flux densities in the left, right and center legs shown in \eqref{eq:5}-\eqref{eq:7}, represented as $B_{\mathrm{o}_1}$, $B_{\mathrm{o}_2}$ and $B_{\mathrm{c}}$ respectively:
\begin{align}
B_{\mathrm{o}_1} &= B'_{\mathrm{o}_1} + B''_{\mathrm{o}_1} \label{eq:5}\\
B_{\mathrm{o}_2} &= B'_{\mathrm{o}_2} - B''_{\mathrm{o}_2} \label{eq:6}\\
B_{\mathrm{c}} &= B''_{\mathrm{c}} \label{eq:7}
\end{align}

\subsubsection*{\textup{Case 1:  $D \le 0.5 \frac{n V_{\mathrm{o}}}{V_{\mathrm{in}_{\mathrm{max}}}}$}}
The waveforms in Fig. 4 represent the circuit operation for Case 1. As shown, flux cancellation does not occur in the outer legs; however, the magnitude of the peak flux density remains uniform across all legs. If the magnetizing current of one phase reaches zero before the magnetizing current of the other phase reaches its peak, then $B_{\mathrm{max}} =  B'_{\mathrm{max}}$. This condition is satisfied when $D$ is expressed by \eqref{eq:8}.
\begin{equation}
\label{eq:8}
D \le 0.5 \frac{n V_{\mathrm{o}}}{V_{\mathrm{in}_{\mathrm{max}}}}
\end{equation}

\subsubsection*{\textup{Case 2: $0.5 \frac{n V_{\mathrm{o}}}{V_{\mathrm{in}_{\mathrm{max}}}} < D \le 0.5$}}
In this case, the magnetizing current of one phase does not reach zero before the subsequent phase reaches its peak. The saturation flux density across all phases can be derived as follows. At $t_2$, $B'_{\mathrm{o}_1} = B'_{\mathrm{max}}$. Moreover, $B''_{\mathrm{o}_1}(t_2) = B'_{\mathrm{o}_1}(t_5)$ and is expressed in \eqref{eq:9}:
\begin{equation}
\label{eq:9}
B''_{\mathrm{o}_1}(t_2) = B'_{\mathrm{o}_1}(t_5) = B'_{\mathrm{max}} - \frac{B'_{\mathrm{max}}}{t_6 - t_2}(t_5 - t_2)
\end{equation}
with $t_5 = 0.5 T_{\mathrm{s}} + t_2$ and $t_2 = D T_{\mathrm{s}}$. The peak $i_{L_{\mathrm{m}_{1_{\mathrm{pk}}}}}$ can be rewritten as in \eqref{eq:10}:
\begin{equation}
\label{eq:10}
i_{L_{{{\mathrm{m}_1}_{\mathrm{pk}}}}} = \frac{V_{\mathrm{in}_{\mathrm{max}}}}{L_{\mathrm{m}}} D T_{\mathrm{s}} = \frac{n V_{\mathrm{o}}}{L_{\mathrm{m}}}(t_6 - t_2)
\end{equation}

In this case, \eqref{eq:8} is not satisfied, and this results in $B_{\mathrm{max}}$ in the left outer leg exceeding $B'_{\mathrm{max}}$ in the center leg, as expressed in \eqref{eq:11}:
\begin{equation}
\label{eq:11}
B_{\mathrm{max}} = \left(2 - \frac{0.5 n V_{\mathrm{o}}}{V_{\mathrm{in}_{\mathrm{max}}} D}\right) B'_{\mathrm{max}}
\end{equation}

\subsubsection*{\textup{Case 3: $0.5 < D$}}
In this case, $B_{\mathrm{max}}$ in the left outer leg consistently exceeds $B'_{\mathrm{max}}$. By analysis similar to Case 2, $B_{\mathrm{max}}$ is derived as \eqref{eq:12}:
\begin{equation}
\label{eq:12}
B_{\mathrm{max}} = \left(1 + \frac{D - 0.5}{D}\right) B'_{\mathrm{max}}
\end{equation}

Because the flyback is interleaved, it is expected that the rectifier will primarily operate at $D \le 0.5$, which operates in the conditions set by Case 1 or 2.

\subsection{Design Method of the Proposed Integrated Transformer}
\begin{figure}[htbp]
\centering
% (a) Top circuit
        \subfloat[]{
        \includegraphics[width=0.6\columnwidth]{Figures_v2/Fig8a.pdf} % Make sure this file exists
        \label{fig:8a}
        }
        \\ % This double backslash forces the next image to the line below (vertical stack)
        % (b) Middle circuit
        \subfloat[]{
        \includegraphics[width=0.75\columnwidth]{Figures_v2/Fig8b.pdf}
        \label{fig:8b}
        }
        
        \caption{EE core configuration: (a) magnetic resistances of the transformer ($R_{\mathrm{mo}}$ for the outer legs, $R_{\mathrm{mc}}$ for the center leg), and (b) effective cross-sectional area ($A_{\mathrm{mo}}$ for the outer legs, $A_{\mathrm{mc}}$ for the center leg) and effective magnetic path length $l_{\mathrm{mo}}$ for the outer legs, $l_{\mathrm{mc}}$ for the center leg.}
\end{figure}
\begin{figure}[htbp]
\centering
% (a) Top circuit
        \subfloat[]{
        \includegraphics[width=0.7\columnwidth]{Figures_v2/Fig9b.pdf} % Make sure this file exists
        \label{fig:9a}
        }
        \\ % This double backslash forces the next image to the line below (vertical stack)
        % (b) Middle circuit
        \subfloat[]{
        \includegraphics[width=0.7\columnwidth]{Figures_v2/Fig9a.pdf}
        \label{fig:9b}
        }
        
        \caption{EE core configuration with air gaps: (a) setup when $l_{\mathrm{mo}_{\mathrm{air}}} > l_{\mathrm{mc}_{\mathrm{air}}}$ and (b) setup when $l_{\mathrm{mo}_{\mathrm{air}}} < l_{\mathrm{mc}_{\mathrm{air}}}$.}
\end{figure}

A turns ratio $n$ is selected in consideration of the maximum voltage ratings of the switches. The maximum operating duty cycle $D$ is selected. The magnetizing inductance, $L_{\mathrm{m}_1}=L_{\mathrm{m}_2}=L_{\mathrm{m}}$ is calculated from the circuit specifications. The saturation magnetic flux density per phase $B'_{\mathrm{max}}$ is calculated from the $B_{\mathrm{max}}$ given by the chosen core material, considering the maximum $D$. A core is chosen using the area-product (AP) method. 

The number of turns is calculated using $B'_{\mathrm{max}}$ and the peak primary side current $I_{\mathrm{pri}_{\mathrm{pk}}}$ as shown in \eqref{eq:13} and \eqref{eq:14}:
\begin{equation}
\label{eq:13}
\frac{N_{12}}{n} = N_{34} \ge \frac{L_{\mathrm{m}} I_{\mathrm{pri}_{\mathrm{pk}}}}{n A_{\mathrm{mo}} B'_{\mathrm{max}}}
\end{equation}
\begin{equation}
\label{eq:14}
\frac{N_{56}}{n} = N_{78} \ge \frac{L_{\mathrm{m}} I_{\mathrm{pri}_{\mathrm{pk}}}}{n A_{\mathrm{mc}} B'_{\mathrm{max}}}
\end{equation}

The magnetic resistances of the outer legs $R_{\mathrm{mo}}$ and the center leg $R_{\mathrm{mc}}$ are obtained from a magnetic circuit analysis of the integrated transformer in its separat phases. The transformer $N_{12}-N_{34}$ is only made of the outer legs as shown in Fig. \ref{fig:8b}, while $N_{56}-N_{78}$ is made of all the legs. The magnetizing inductance $L_{\mathrm{m}}$ in terms of the number of turns and magnetic resistances for $N_{12}-N_{34}$ and $N_{56}-N_{78}$ are shown in \eqref{eq:15} and \eqref{eq:16} respectively.
\begin{align}
L_{\mathrm{m}} &= \frac{N_{12}^2}{2 R_{\mathrm{mo}}} \label{eq:15}\\
L_{\mathrm{m}} &= \frac{N_{56}^2}{R_{\mathrm{mc}} + 0.5 R_{\mathrm{mo}}} \label{eq:16}
\end{align}

The magnetic resistances $R_{\mathrm{mo}}$ and $R_{\mathrm{mc}}$ are obtained from \eqref{eq:17} and \eqref{eq:18}:
\begin{align}
R_{\mathrm{mo}} &= \frac{N_{12}^2}{2 L_{\mathrm{m}}} \label{eq:17}\\
R_{\mathrm{mc}} &= \frac{(2 N_{56} - N_{12})(2 N_{56} + N_{12})}{4 L_{\mathrm{m}}} \label{eq:18}
\end{align}

Finally, an air gap was also inserted in each of the legs to provide the desired resistance. Non-equal length air gaps were placed in the transformer. There are two cases, one where the air gaps in the outer legs $l_{\mathrm{mo}_{\mathrm{air}}}$ are longer than the air gap in the center leg $l_{\mathrm{mc}_{\mathrm{air}}}$, and one where $l_{\mathrm{mc}_{\mathrm{air}}}$ is longer than $l_{\mathrm{mo}_{\mathrm{air}}}$.

\subsubsection*{\textup{Case 1: $l_{\mathrm{mo}_{\mathrm{air}}} > l_{\mathrm{mc}_{\mathrm{air}}}$}}
Given the setup in Fig. \ref{fig:9a}, $R_{\mathrm{mo}}$ and $R_{\mathrm{mc}}$ are shown in \eqref{eq:19} and \eqref{eq:20} respectively with $l_{\mathrm{mo}_{\mathrm{air}}}$ and $l_{\mathrm{mc}_{\mathrm{air}}}$ in \eqref{eq:21} and \eqref{eq:22} respectively.  
\begin{align}
R_{\mathrm{mo}} &= \frac{l_{\mathrm{mo}} - l_{\mathrm{mo}_{\mathrm{air}}} + l_{\mathrm{mc}_{\mathrm{air}}}}{\mu_r \mu_0 A_{\mathrm{mo}}} + \frac{l_{\mathrm{mo}_{\mathrm{air}}}}{\mu_0 A_{\mathrm{mo}}} \label{eq:19}\\
R_{\mathrm{mc}} &= \frac{l_{\mathrm{mc}}}{\mu_r \mu_0 A_{\mathrm{mc}}} + \frac{l_{\mathrm{mc}_{\mathrm{air}}}}{\mu_0 A_{\mathrm{mc}}} \label{eq:20}\\
l_{\mathrm{mo}_{\mathrm{air}}} &= \frac{\mu_0 (\mu_r A_{\mathrm{mo}} R_{\mathrm{mo}} - A_{\mathrm{mc}} R_{\mathrm{mc}}) + \frac{l_{\mathrm{mc}}}{\mu_r} - l_{\mathrm{mo}}}{\mu_r - 1} \label{eq:21}\\
l_{\mathrm{mc}_{\mathrm{air}}} &= \mu_0 A_{\mathrm{mc}} R_{\mathrm{mc}} - \frac{l_{\mathrm{mc}}}{\mu_r} \label{eq:22}
\end{align}

\subsubsection*{\textup{Case 2: $l_{\mathrm{mo}_{\mathrm{air}}} < l_{\mathrm{mc}_{\mathrm{air}}}$}}
Given the setup in Fig. \ref{fig:9b}, $R_{\mathrm{mo}}$ and $R_{\mathrm{mc}}$ are shown in \eqref{eq:23} and \eqref{eq:24} respectievly $l_{\mathrm{mo}_{\mathrm{air}}}$ and $l_{\mathrm{mc}_{\mathrm{air}}}$ in \eqref{eq:25} and \eqref{eq:26} respectively.
\begin{align}
R_{\mathrm{mo}} &= \frac{l_{\mathrm{mo}}}{\mu_r \mu_0 A_{\mathrm{mo}}} + \frac{l_{\mathrm{mo}_{\mathrm{air}}}}{\mu_0 A_{\mathrm{mo}}} \label{eq:23}\\
R_{\mathrm{mc}} &= \frac{l_{\mathrm{mc}} - l_{\mathrm{mc}_{\mathrm{air}}} + l_{\mathrm{mo}_{\mathrm{air}}}}{\mu_r \mu_0 A_{\mathrm{mc}}} + \frac{l_{\mathrm{mc}_{\mathrm{air}}}}{\mu_0 A_{\mathrm{mc}}} \label{eq:24}\\
l_{\mathrm{mo}_{\mathrm{air}}} &= \mu_0 A_{\mathrm{mo}} R_{\mathrm{mo}} - \frac{l_{\mathrm{mo}}}{\mu_r} \label{eq:25}\\
l_{\mathrm{mc}_{\mathrm{air}}} &= \frac{\mu_0 (\mu_r A_{\mathrm{mc}} R_{\mathrm{mc}} - A_{\mathrm{mo}} R_{\mathrm{mo}}) + \frac{l_{\mathrm{mo}}}{\mu_r} - l_{\mathrm{mc}}}{\mu_r - 1} \label{eq:26}
\end{align}


\section{Simulation Verification of Proposed Design}

\subsection{Simulation Verification of the Proposed Design Method}
To ensure the valid operation of the proposed integrated transformer, a simulation of the two-phase interleaved flyback rectifier was first performed in PLECS Standalone 4.9.4, with the diagram shown in Fig. \ref{fig:10a}. Clamped RCD Snubbers ($C_{\mathrm{sn}_1}, D_{\mathrm{sn}_1}, R_{\mathrm{sn}_1}$ and $C_{\mathrm{sn}_2}, D_{\mathrm{sn}_2}, R_{\mathrm{sn}_2}$) were also added for protection against surge voltage \cite{ref15, ref16} caused by the leakage inductance and the switch capacitance at the turn-off of $S_1$ and $S_2$. An integrated transformer, shown in Fig. 10(b), was also designed using the proposed method. The circuit and transformer were designed with the condition given by \eqref{eq:8}, and will thus be verified in this study. The specifications of the two-phase interleaved flyback rectifier and transformer for the simulation verification are shown in Tables \ref{tab:sim_rectifier} and \ref{tab:sim_transformer} respectively. The duty cycle is calculated using \eqref{eq:27} \cite{ref8}.
\begin{equation}
\label{eq:27}
D = \frac{2}{V_{\mathrm{in}_{\mathrm{pk}}}} \sqrt{\frac{0.5 L_{\mathrm{m}} P_{\mathrm{o}}}{T_{\mathrm{s}}}}
\end{equation}

\begin{table}[b!]
    \caption{Simulation Specifications of the Prototype Two-Phase Interleaved Flyback Rectifier}
    \centering
    \footnotesize
    \setlength{\tabcolsep}{3pt} % Slightly tighter spacing to handle long text
    \begin{tabular}{ccc}
        \hline
        \textbf{Item} & \textbf{Symbol} & \textbf{Value} \\
        \hline \hline
        AC input voltage & $v_{\mathrm{in}}$ & 100 V V$_{\textrm{rms}}$\\
        Maximum rated power & $P_{\mathrm{O}}$ & 240 W \\
        Output voltage & $V_{\mathrm{O}}$ & 48 V \\
        Switching frequency & $f_{\mathrm{s}}$ & 100 kHz \\
        Input capacitor & $C_{\mathrm{in}}$ & 1 \textmu F \\
        Output capacitor & $C_{\mathrm{O}}$ & 2 x 1mF (parallel) \\
        Snubber capacitor & $C_{\mathrm{sn}_1}, C_{\mathrm{sn}_2}$ & 2.2 nF \\
        Snubber resistor & $R_{\mathrm{sn}_1}, R_{\mathrm{sn}_2}$ & 30 k$\Omega$ \\
        \hline
    \end{tabular}
    \label{tab:sim_rectifier}
\end{table}

\begin{table}[b!]
    \caption{Simulation Specifications of the Prototype Integrated Transformer}
    \centering
    \footnotesize
    \setlength{\tabcolsep}{4pt} 
    \begin{tabular}{ccc}
        \hline
        \textbf{Item} & \textbf{Symbol} & \textbf{Value} \\
        \hline \hline
        Magnetic core & - & OP44721EC \\
        Relative permeability & $\mu_r$ & 2500\\
        Magnetic path length (center leg) & $l_{\mathrm{mc}}$ & 31.7 mm\\
        Magnetic path length (outer leg) & $l_{\mathrm{mo}}$ & 71.4 mm\\
        Cross-sectional area (center leg) & $A_{\mathrm{mc}}$ & 243.4 mm$^2$ \\
        Cross-sectional area (outer leg) & $A_{\mathrm{mo}}$ & 117.6 mm$^2$\\
        Magnetizing inductance & $L_{\mathrm{m}_1}, L_{\mathrm{m}_2}$ & 82.8 \textmu H \\
        Turn ratio & $n$ & 3 \\
        \multirow{2}{*}{Number of turns} & $N_{12}:N_{34}$ & 24:8 \\
                                         & $N_{56}:N_{78}$ & 18:6 \\
        Center leg air gap & $l_{{\mathrm{mc}}_{\mathrm{air}}}$ & 0.650 mm \\
        Outer leg air gap & $l_{{\mathrm{mo}}_{\mathrm{air}}}$ & 0.500 mm \\
        Coupling coefficient of $N_{12}, N_{34}$ & $k_{N_{12}-N_{34}}$ & 0.990 \\
        Coupling coefficient of $N_{56}, N_{78}$ & $k_{N_{56}-N_{78}}$ & 0.990 \\
        Coupling coefficient of $N_{12}, N_{56}$ & $k_{N_{12}-N_{56}}$ & 0.000 \\
        \hline
    \end{tabular}
    \label{tab:sim_transformer}
\end{table}

\begin{figure}[htbp]
\centering
% (a) Top circuit
        \subfloat[]{
        \includegraphics[width=\columnwidth]{Figures_v2/Fig10a.pdf} % Make sure this file exists
        \label{fig:10a}
        }
        \\ % This double backslash forces the next image to the line below (vertical stack)
        % (b) Middle circuit
        \subfloat[]{
        \includegraphics[width=0.6\columnwidth]{Figures_v2/Fig10b.pdf}
        \label{fig:10b}
        }
        
        \caption{Prototype two-phase interleaved flyback rectifier (a) with the proposed integrated transformer (b).}
\end{figure}

\subsection{Simulation Results}
Fig. \ref{fig:11} shows waveforms from the simulation testing of the two-phase interleaved flyback rectifier with the proposed transformer at 240 W. Key waveforms of the circuit were observed from the simulation, which include $v_{\mathrm{gs}_1}$ and $v_{\mathrm{gs}_2}$, $i_{L_{\mathrm{m}_1}}$ and $i_{L_{\mathrm{m}_2}}$, and the fluxes $B_{\mathrm{o}_1}$, $B_{\mathrm{o}_2}$, and $B_{\mathrm{c}}$. Similar waveforms to that of Fig. \ref{fig:4} have been observed, confirming that the system operates exactly as expected from the theoretical analysis. Flux cancellation is observed in the center leg during the time that $N_{56}-N_{78}$ is inactive, as denoted by the interval when $i_{L_{\mathrm{m}_2}} = 0$. The magnetizing currents $i_{L_{\mathrm{m}_1}}$ and $i_{L_{\mathrm{m}_2}}$ reach zero before one other phase reaches its peak. Moreover, the increase and decrease of fluxes were all linear, therefore confirming that there is indeed no magnetic coupling interference. This confirms the validity of the theoretical operation proposed in Section II.

\begin{figure}[b!]
    % Import the separate file here
    \includegraphics[width=\columnwidth]{Figures_v2/sim_plot_output.pdf}
    \caption{Key simulation waveforms of the proposed two-phase interleaved flyback rectifier at 240 W.}
    \label{fig:11}
\end{figure}

\section{Size Reduction Effect of the Proposed Transformer}
This section examines the size reduction effect by the integrated transformer. For this study, the conventional transformer is assumed to be two separate transformers that facilitate isolation between the primary and secondary side. A representation is shown in Fig. \ref{fig:12a}, and its windings are represented in Fig. \ref{fig:13a}. The proposed integrated transformer is represented in Fig. \ref{fig:12b} and its windings are represented  in Fig. \ref{fig:13b}. The volume is compared by using the AP method. The AP method utilizes the product between the core cross-sectional area, $A_{\mathrm{c}}$, and the window-area, $A_{\mathrm{w}}$, and is a fundamental core size for the application \cite{ref21}. It can be expressed both in terms of the transformer geometric parameters and the circuit parameters, as will be demonstrated in the following analysis.
\begin{figure}[b!]
\centering
% (a) Top circuit
        \subfloat[]{
        \includegraphics[width=0.9\columnwidth]{Figures_v2/Fig12a.pdf} % Make sure this file exists
        \label{fig:12a}}
        \\ % This double backslash forces the next image to the line below (vertical stack)
        % (b) Middle circuit
        \subfloat[]{
        \includegraphics[width=0.7\columnwidth]{Figures_v2/Fig12b.pdf}
        \label{fig:12b}}
        
        \caption{Two separate transformers (a) and proposed integrated transformer (b) used in the volume analysis.}
        \label{fig:12}
\end{figure}
\begin{figure}[htbp]
\centering
% (a) Top circuit
        \subfloat[]{
        \includegraphics[width=0.9\columnwidth]{Figures_v2/Fig13a.pdf} % Make sure this file exists
        \label{fig:13a}
        }
        \\ % This double backslash forces the next image to the line below (vertical stack)
        % (b) Middle circuit
        \subfloat[]{
        \includegraphics[width=0.7\columnwidth]{Figures_v2/Fig13b.pdf}
        \label{fig:13b}
        }
        
        \caption{Two separate transformers (a) and the proposed integrated transformer (b) used in the volume analysis, with the windings represented as well.}
        \label{fig:13}
\end{figure}

The EE-core was redesigned in such a way such that $A_{\mathrm{w}}$ and $A_{\mathrm{c}}$ are squares to allow for simpler analysis, as shown in Fig. \ref{fig:14} \cite{ref22}. Here, $A_{\mathrm{c}}$ is twice the area of the outer leg.

\begin{figure}[htbp]
\centering
        \includegraphics[width=\columnwidth]{Figures_v2/Fig14.pdf} % Make sure this file exists
        
        \caption{Simplified representation of one EE-core used in the volume analysis.}
        \label{fig:14}
\end{figure}


For the conventional transformer, $A_{\mathrm{c}}$ is expressed in \eqref{eq:28}. Here, $N_{\mathrm{p}}$ is the number of turns wound on the center leg of the conventional transformer.
\begin{equation}
\label{eq:28}
A_{\mathrm{c}_{\mathrm{conv}}} = \frac{L_{\mathrm{m}} I_{\mathrm{pri}_{\mathrm{pk}}}}{B_{\mathrm{max}} N_{\mathrm{p}}}
\end{equation}
For the integrated transformer $A_{\mathrm{c}}$ is expressed in \eqref{eq:29}.
\begin{equation}
\label{eq:29}
A_{\mathrm{c}_{\mathrm{integ}}} = \frac{L_{\mathrm{m}} I_{\mathrm{pri}_{\mathrm{pk}}}}{B_{\mathrm{max}} N_{56}}
\end{equation}

The window area $A_{\mathrm{w}}$ can be expressed in terms of winding area and the occupancy ratio $K_{\mathrm{u}}$, as seen in \eqref{eq:30} for the conventional transformer:
\begin{equation}
\label{eq:30}
A_{\mathrm{w}_{\mathrm{conv}}} K_{\mathrm{u}} = N_{\mathrm{p}} \frac{I_{\mathrm{pri}_{\mathrm{rms}}}}{J_{\mathrm{m}}} + N_{\mathrm{s}} \frac{I_{\mathrm{sec}_{\mathrm{rms}}}}{J_{\mathrm{m}}}
\end{equation}
which simplifies to:
\begin{equation}
\label{eq:30b}
A_{\mathrm{w}_{\mathrm{conv}}} = \frac{N_{\mathrm{p}}}{J_{\mathrm{m}} K_{\mathrm{u}}} \left(I_{\mathrm{pri}_{\mathrm{rms}}} + \frac{I_{\mathrm{sec}_{\mathrm{rms}}}}{n}\right)
\end{equation}
For the integrated transformer, taking into account that the core window has half the total turns of the transformer $N_{12}$ and $N_{34}$, the $A_{\mathrm{w}}$ can be expressed as seen in \eqref{eq:32}.
\begin{align}
\label{eq:32}
A_{\mathrm{w}_{\mathrm{integ}}} K_{\mathrm{u}} &= \left(N_{56} \frac{I_{\mathrm{pri}_{\mathrm{rms}}}}{J_{\mathrm{m}}} + N_{78} \frac{I_{\mathrm{sec}_{\mathrm{rms}}}}{J_{\mathrm{m}}}\right) \nonumber \\
&\quad + \frac{1}{2}\left(N_{12} \frac{I_{\mathrm{pri}_{\mathrm{rms}}}}{J_{\mathrm{m}}} + N_{34} \frac{I_{\mathrm{sec}_{\mathrm{rms}}}}{J_{\mathrm{m}}}\right)
\end{align}

The total volume of the EE core in Fig. \ref{fig:14} expressed in \eqref{eq:33}, where $x$ is a constant that denotes the ratio $A_{\mathrm{c}}/A_{\mathrm{w}}$. The volume corresponding to the conventional and integrated transformer is obtained using their respective $A_{\mathrm{c}}$ and $A_{\mathrm{w}}$.
\begin{equation}
\label{eq:33}
\mathrm{Vol} = 2 AP^{\frac{3}{4}} (2 x^{\frac{1}{4}} + x^{\frac{3}{4}} + x^{-\frac{1}{4}})
\end{equation}
The rms expressions of the primary and secondary side current are given in \eqref{eq:34} and \eqref{eq:35} \cite{ref8}.
\begin{align}
I_{\mathrm{pri}_{\mathrm{rms}}} &= \frac{V_{\mathrm{in}_{\mathrm{max}}}}{L_{\mathrm{m}} f_{\mathrm{s}}} \sqrt{\frac{D^3}{6}} \label{eq:34}\\
I_{\mathrm{sec}_{\mathrm{rms}}} &= \frac{V_{\mathrm{in}_{\mathrm{max}}}}{L_{\mathrm{m}} f_{\mathrm{s}}} \sqrt{\frac{4 n D^3}{9 \pi} \cdot \frac{V_{\mathrm{in}_{\mathrm{max}}}}{V_{\mathrm{o}}}} \label{eq:35}
\end{align}

Using these equations with the actual circuit specfications, a 56\% volume reduction has been achieved by the proposed design, with the comparison shown in Table \ref{tab:volume}. This assumes that the integrated core geometry is custom-tailored to the specific area product requirements derived from the circuit specifications, rather than constrained by standard commercial core shapes. While the experimental prototype described in Section V utilizes a commercial core (Magnetics 0P44721EC) for proof-of-concept, this theoretical analysis confirms that a purpose-built integrated core would offer significant improvements in power density compared to the conventional discrete implementation.

\begin{table}[b!]
    \caption{Volume Analysis of the Conventional and Proposed Transformer}
    \centering
    \footnotesize
    \setlength{\tabcolsep}{2pt}
    \begin{tabular}{cccc}
        \hline
        \textbf{Item} & \textbf{Symbol} & \textbf{Conventional} & \textbf{Proposed} \\
        \hline \hline
        Maximum input voltage & $V_{\mathrm{in}_{\mathrm{max}}}$ & \multicolumn{2}{c}{240 V V$_{\mathrm{rms}}$} \\
        Output voltage & $V_{\mathrm{in}_{\mathrm{max}}}$ & \multicolumn{2}{c}{48 V} \\
        Maximum duty cycle & $D$ & \multicolumn{2}{c}{0.47} \\
        Switching frequency & $f_\mathrm{s}$ & \multicolumn{2}{c}{100 kHz} \\
        \hline
        Magnetizing inductance & $L_{\mathrm{m}}$ & \multicolumn{2}{c}{82.8 \textmu H} \\
        Turn ratio & $n$ & \multicolumn{2}{c}{3} \\
        
        % --- TURNS SECTION ---
        Number of turns (Conv.) & $N_{\mathrm{p}}:N_{\mathrm{s}}$ & 24:8 & \multicolumn{1}{c}{--} \\
        
        % Multirow for Integrated
        \multirow{2}{*}{Number of turns (Prop.)} & $N_{12}:N_{34}$ & \multicolumn{1}{c}{--} & 24:8 \\
                                                      & $N_{56}:N_{78}$ & \multicolumn{1}{c}{--} & 18:6 \\
        
        % --- SHARED PARAMETERS ---
        Primary side peak current & $I_{\mathrm{pri,pk}}$ & \multicolumn{2}{c}{8.02 A} \\
        Primary side rms current & $I_{\mathrm{pri,rms}}$ & \multicolumn{2}{c}{5.39 A} \\
        Secondary side rms current & $I_{\mathrm{sec,rms}}$ & \multicolumn{2}{c}{12.53 A} \\
        Saturation flux density & $B_{\mathrm{max}}$ & \multicolumn{2}{c}{0.3 T} \\
        Maximum current density & $J_{\mathrm{m}}$ & \multicolumn{2}{c}{5 A/mm$^2$} \\
        Occupancy ratio & $K_{\mathrm{u}}$ & \multicolumn{2}{c}{0.3} \\
        \hline
        Area Product & $\mathrm{AP}$ & 14128.75 mm$^4$ & 23547.91 mm$^4$ \\
        Core Volume & $\mathrm{Vol}$ & 21106.38 mm$^3$ & 9398.46 mm$^3$ \\
        \hline
    \end{tabular}
    \label{tab:volume}
\end{table}

\section{Experimental Verification of Proposed Design}

\subsection{Prototype Evaluation Circuit and Prototype Transformer}
A prototype two-phase interleaved flyback rectifier, shown in Fig. \ref{fig:10a} operating in DCM was built to verify the operation of the integrated transformer. The prototype integrated transformer, shown in Fig. \ref{fig:10b}, was also designed using the proposed method. The specifications of the two-phase interleaved flyback rectifier and transformer are shown in Tables \ref{tab:prototype_rectifier} and \ref{tab:prototype_transformer} respectively. The transformer parameters were characterized using a Hioki IM3536 LCR meter at 100 kHz. The coupling coefficients of the transformers $k_{N_{12}-N_{34}}$, $k_{N_{56}-N_{78}}$, and $k_{N_{12}-N_{56}}$ were derived from the mutual inductance obtained after connecting the respective transformer windings in series-aiding and series-opposing configurations. The coupling coefficient between the two flyback transformers $k_{N_{12}-N_{56}}$ is 0.0005, showing that the coupling was almost eliminated.



\subsection{Experimental Results}
Fig. \ref{fig:16a} shows the 100 W experimental waveforms at 4 ms/div from the two-phase interleaved flyback rectifier with the proposed transformer, demonstrating power factor correction. Fig. \ref{fig:16b} presents 2 \textmu s/div waveforms. No current is induced in either switch except during the resonance between the leakage inductance and the parasitic capacitance of the switch. During the inactive period of the transformer, secondary diode side current is zero. Moreover, after $i_{\mathrm{S}_1}$ and $i_{\mathrm{S}_2}$ reached zero when $S_1$ and $S_2$ switched off, the currents remained at zero. Therefore, the two flyback converters operate independently without magnetic coupling interference. The same phenomenon was observed at 240 W two-phase drive, as shown in Fig. \ref{fig:16c}. Single-phase drive at lower power at 32.5 W is also feasible as shown in Fig. \ref{fig:16d}.

% --- TABLE VI: Rectifier Specifications ---
\begin{table}[t!]
    \caption{Specifications of the Prototype Two-Phase Interleaved Flyback Rectifier}
    \centering
    \footnotesize
    \setlength{\tabcolsep}{3pt} % Slightly tighter spacing to handle long text
    \begin{tabular}{ccc}
        \hline
        \textbf{Item} & \textbf{Symbol} & \textbf{Value / Component} \\
        \hline \hline
        AC input voltage & $v_{\mathrm{in}}$ & 100-240 V  V$_{\textrm{rms}}$ \\
        Maximum rated power & $P_{\mathrm{O}}$ & 240 W \\
        Output voltage & $V_{\mathrm{O}}$ & 48 V \\
        Switching frequency & $f_{\mathrm{s}}$ & 100 kHz \\
        Bridge diode & $BD$ & CD-DF408S \\
        Input capacitor & $C_{\mathrm{in}}$ & 1 \textmu F \\
        Switch & $S_1, S_2$ & STF15N80k5 \\
        Secondary diode & $D_1, D_2$ & STPSC30H12CWL \\
        Output capacitor & $C_{\mathrm{O}}$ & 2 x 1mF, in parallel \\
        Snubber capacitor & $C_{\mathrm{sn}_1}, C_{\mathrm{sn}_2}$ & 2.2 nF \\
        Snubber resistor & $R_{\mathrm{sn}_1}, R_{\mathrm{sn}_2}$ & 30 k$\Omega$ \\
        Snubber diode & $D_{\mathrm{sn}_1}, D_{\mathrm{sn}_2}$ & VS-E7MH0112-M3/1 \\
        \hline
    \end{tabular}
    \label{tab:prototype_rectifier}
\end{table}

% --- TABLE VII: Transformer Specifications ---
\begin{table}[t!]
    \caption{Specifications of the Prototype Integrated Transformer}
    \centering
    \footnotesize
    \setlength{\tabcolsep}{4pt} 
    \begin{tabular}{ccc}
        \hline
        \textbf{Item} & \textbf{Symbol} & \textbf{Value} \\
        \hline \hline
        Magnetizing inductance (Target) & $L_{\mathrm{m}}$ & 82.8 \textmu H \\
        Magnetizing inductance of $N_{12}-N_{34}$ & $L_{\mathrm{m}_1}$ & 86.1 \textmu H \\
        Magnetizing inductance of $N_{56}-N_{78}$ & $L_{\mathrm{m}_2}$ & 75.1 \textmu H \\
        Turn ratio & $n$ & 3 \\
        Magnetic core & - & OP44721EC \\
        \multirow{2}{*}{Number of turns} & $N_{12}:N_{34}$ & 24:8 \\
                                         & $N_{56}:N_{78}$ & 15:5 \\
        Center leg air gap & $l_{{\mathrm{mc}}_{\mathrm{air}}}$ & 0.650 mm \\
        Outer leg air gap & $l_{{\mathrm{mo}}_{\mathrm{air}}}$ & 0.500 mm \\
        Coupling coefficient of $N_{12}, N_{34}$ & $k_{N_{12}-N_{34}}$ & 0.970 \\
        Coupling coefficient of $N_{56}, N_{78}$ & $k_{N_{56}-N_{78}}$ & 0.979 \\
        Coupling coefficient of $N_{12}, N_{56}$ & $k_{N_{12}-N_{56}}$ & 0.0005 \\
        \hline
    \end{tabular}
    \label{tab:prototype_transformer}
\end{table}

\begin{figure}[htbp]
\centering
% (a) Top circuit
        \subfloat[]{
        \includegraphics[width=0.9\columnwidth]{Figures_v2/Fig16a.pdf} % Make sure this file exists
        \label{fig:16a}}
        \\ % This double backslash forces the next image to the line below (vertical stack)
        % (b) Middle circuit
        \subfloat[]{
        \includegraphics[width=0.9\columnwidth]{Figures_v2/Fig16b.pdf}
        \label{fig:16b}}
        \\ % This double backslash forces the next image to the line below (vertical stack)
        % (b) Middle circuit
        \subfloat[]{
        \includegraphics[width=0.9\columnwidth]{Figures_v2/Fig16c.pdf}
        \label{fig:16c}}
        \\ % This double backslash forces the next image to the line below (vertical stack)
        % (b) Middle circuit
        \subfloat[]{
        \includegraphics[width=0.9\columnwidth]{Figures_v2/Fig16d.pdf}
        \label{fig:16d}}
        
        \caption{Experimental waveforms of a two-phase interleaved flyback rectifier with the proposed integrated transformer at (a) \& (b): 100 W, two-phase AC drive, (c): 240 W, two-phase AC drive, (d): 32.5 W, single-phase AC drive. For (a) and (b): $v_{\mathrm{in}}$: 50 V/div., $V_{\mathrm{o}}$: 10 V/div., $v_{\mathrm{ds}_1}, v_{\mathrm{ds}_2}$: 100 V/div., $i_{\mathrm{S}_1}, i_{\mathrm{S}_2}$: 1 A/div., $i_{\mathrm{D}_1}, i_{\mathrm{D}_2}$: 5 A/div.; For (c): $v_{\mathrm{in}}$: 50 V/div., $V_{\mathrm{o}}$: 10 V/div., $v_{\mathrm{ds}1}, v_{\mathrm{ds}2}$: 100 V/div., $i_{\mathrm{S}_1}, i_{\mathrm{S}_2}$: 2 A/div., $i_{\mathrm{D}_1}, i_{\mathrm{D}_2}$: 5 A/div.; For (d): $v_{\mathrm{in}}$: 50 V/div., $V_{\mathrm{o}}$: 20 V/div., $v_{\mathrm{ds}_1}, v_{\mathrm{ds}_2}$: 100 V/div., $i_{\mathrm{S}_1}, i_{\mathrm{S}_2}$: 2 A/div., $i_{\mathrm{D}_1}, i_{\mathrm{D}_2}$: 2 A/div.}
        \label{fig:16}
\end{figure}

\section{Conclusion}
In this paper, an integrated transformer was successfully designed for a two-phase interleaved flyback rectifier for USB PD 3.1 applications. The proposed winding structure for the integrated transformer prevents coupling between the two flyback transformers and achieves simple manufacturability through the use of a commercial EE-core. A coupling coefficient of 0.0005 was achieved with the prototype. Independent flyback circuit operation without magnetic interference due to coupling was verified through simulation and experiment up to 240 W. Because of the proposed design, two transformers are integrated into one component, enabling the downsizing of magnetic components by about 56\% compared to the conventional transformer.

% --- TABLE V: Volume Analysis ---




% --- TABLE VII: Transformer Specifications ---
% --- TABLE VII: Transformer Specifications ---


\begin{thebibliography}{1}
\bibliographystyle{IEEEtran}

\bibitem{ref1}
A. Medina-Garcia, J. Cruz-Cozar, D. P. Morales-Santos, N. Rodriguez, and M. Schlenk, "A high-efficiency isolated wide voltage range DC-DC converter using WBG devices," \textit{IEEE Access}, vol. 10, pp. 84818-84825, Aug. 2022.

\bibitem{ref2}
USB-IF PR, "USB Promoter Group announces USB Power Delivery specification revision 3.1," 2021. [Online]. Available: https://www.usb.org/sites/default/files/2021-05/USB\%20PG\%20USB\%20PD\%203.1\%20DevUpdate\%20\\Announcement\_FINAL.pdf

\bibitem{ref8}
J. Shin, S. Choi, and B. Cho, "High-efficiency bridgeless flyback rectifier with bidirectional switch and dual output windings," \textit{IEEE Trans. Power Electron.}, vol. 29, no. 9, pp. 4752-4762, Oct. 2013.

\bibitem{ref9}
S. Choi, J. Shin, J. Imaoka, and M. Yamamoto, "Common-mode noise reduction of bridgeless flyback PFC rectifier with balance technique," \textit{IEICE Tech. Rep.}, vol. 123, no. 119, pp. 67-72, 2023.

\bibitem{ref10}
K. T. Mok, Y. M. Lai, and K. H. Loo, "A single-stage bridgeless power-factor-correction rectifier based on flyback topology," in \textit{Proc. IEEE 33rd Int. Telecommun. Energy Conf. (INTELEC)}, 2011, pp. 1-6.

\bibitem{ref11}
J. Beak, J. Shin, P. Jang, and B. Cho, "A critical conduction mode bridgeless flyback converter," in \textit{Proc. 8th Int. Conf. Power Electron. - ECCE Asia}, 2011, pp. 487-492.

\bibitem{ref3}
B. C. Kim, K. B. Park, C. E. Kim, B. H. Lee, and G. W. Moon, "LLC resonant converter with adaptive link-voltage variation for a high-power-density adapter," \textit{IEEE Trans. Power Electron.}, vol. 25, no. 9, pp. 2248-2252, May 2010.

\bibitem{ref4}
Y. C. Lin, F. C. Lee, Q. Li, X. Huang, and Z. Lin, "A novel AC-to-DC adaptor with ultra-high power density and efficiency," in \textit{Proc. IEEE Appl. Power Electron. Conf. Expo. (APEC)}, 2016, pp. 1853-1860.

\bibitem{ref5}
M. Li, Z. Ouyang, and M. A. E. Andersen, "Analysis and optimal design of high-frequency and high-efficiency asymmetrical half-bridge flyback converters," \textit{IEEE Trans. Ind. Electron.}, vol. 67, no. 10, pp. 8312-8321, Nov. 2019.

\bibitem{ref6}
X. Huang, J. Feng, W. Du, F. C. Lee, and Q. Li, "Design consideration of MHz active clamp flyback converter with GaN devices for low power adapter application," in \textit{Proc. IEEE Appl. Power Electron. Conf. Expo. (APEC)}, 2016, pp. 2334-2341.

\bibitem{ref7}
J. Zhang, X. Huang, X. Wu, and Z. Qian, "A high efficiency flyback converter with new active clamp technique," \textit{IEEE Trans. Power Electron.}, vol. 25, no. 7, pp. 1775-1785, Feb. 2010.

\bibitem{ref12}
B. Shafer, "Interleaving contributes unique benefits to forward converters and flyback converters," in \textit{Texas Instruments Power Supply Design Seminar}, Topic, vol. 4, pp. 1-32, 2004.

\bibitem{Su_2013}
J.-T. Su and C.-W. Liu, "A novel phase-shedding control scheme for improved light load efficiency of multiphase interleaved DC-DC converters," \textit{IEEE Trans. Power Electron.}, vol. 28, no. 10, pp. 4742-4752, Oct. 2013.

\bibitem{Zhang_2013}
Y. Zhang, X.-F. He, Z. Zhang, and Y.-F. Liu, "A hybrid control method for photovoltaic grid-connected interleaved flyback micro-inverter to achieve high efficiency in wide load range," in \textit{Proc. IEEE Appl. Power Electron. Conf. Expo. (APEC)}, 2013, pp. 751-756.

\bibitem{Chen_2014}
H.-C. Chen, "Interleaved current sensorless control for multiphase boost-type switch-mode rectifier with phase-shedding operation," \textit{IEEE Trans. Ind. Electron.}, vol. 61, no. 2, pp. 766-775, Feb. 2014.

\bibitem{Yang_2020}
F. Yang, C. Li, Y. Cao, and K. Yao, "Two-phase interleaved boost PFC converter with coupled inductor under single-phase operation," \textit{IEEE Trans. Power Electron.}, vol. 35, no. 1, pp. 169-184, Jan. 2020.

\bibitem{Z_Li_2025}
Z. Li, J. Liang, M. Fu, T. Long, and H. Wang, "Nonlinear coupled inductor-based light-load efficiency boost technique for trans-inductor voltage regulators," \textit{IEEE Trans. Power Electron.}, vol. 40, no. 10, pp. 14368-14373, Oct. 2025.

\bibitem{Ahmad}
U. Ahmad, H. Cha, and N. Naseem, "Integrated current balancing transformer based input-parallel output-parallel LLC resonant converter modules," \textit{IEEE Trans. Power Electron.}, vol. 36, no. 5, pp. 5278-5289, May 2021.

\bibitem{Dou}
Y. Dou, Z. Ouyang, and M. A. E. Andersen, "Integrated coupled inductors with functionality of current balancing transformer for two-phase synchronous DC-DC converters," \textit{IEEE Trans. Power Electron.}, vol. 35, no. 5, pp. 4472-4476, May 2020.

\bibitem{Nakakohara}
Y. Nakakohara, H. Otake, T. M. Evans, T. Yoshida, M. Tsuruya, and K. Nakahara, "Three-phase LLC series resonant DC/DC converter using SiC MOSFETs to realize high-voltage and high-frequency operation," \textit{IEEE Trans. Ind. Electron.}, vol. 63, no. 4, pp. 2103-2110, Apr. 2016.

\bibitem{Z_Chen_s42835}
Z. Chen, L. Shen, and Y. Wu, "An improved interleaved flyback converter with reduced voltage stress and current auto-sharing," \textit{J. Electr. Eng. Technol.}, vol. 19, pp. 2251-2263, Nov. 2023.

\bibitem{Jin_2024}
F. Jin, T. Yuan, A. Nabih, and Q. Li, "Efficient integrated magnetics with winding cancellation technique to reduce common-mode EMI noise for a single-phase CLLC converter," \textit{IEEE Trans. Power Electron.}, vol. 39, no. 11, pp. 14758-14774, Nov. 2024.

\bibitem{Ouyang}
Z. Ouyang, Z. Zhang, M. A. E. Andersen, and O. C. Thomsen, "Four quadrants integrated transformers for dual-input isolated DC-DC converters," \textit{IEEE Trans. Power Electron.}, vol. 27, no. 6, pp. 2697-2702, Jun. 2012.

\bibitem{ref19}
G. B. Koo, "Design guidelines for RCD snubber of flyback converters," Fairchild Semiconductor, San Jose, CA, USA, Appl. Note AN-4147 Rev. 1.1.0, 2006.

\bibitem{ref20}
W. McMurray, "Selection of snubbers and clamps to optimize the design of transistor switching converters," \textit{IEEE Trans. Ind. Appl.}, vol. IA-16, no. 4, pp. 513-523, 1980.

\bibitem{ref21}
W. Colonel and T. McLyman, \textit{Transformer and Inductor Design Handbook}. Boca Raton, FL, USA: CRC Press, 2004.

\bibitem{ref22}
S. Kimura, S. Aoto, M. Ishihara, J. Imaoka, and M. Yamamoto, "Interleaved active clamp forward converter with novel integrated magnetic components," in \textit{Proc. IEEE Energy Convers. Congr. Expo. (ECCE)}, 2015, pp. 6029-6036.

\end{thebibliography}

\begin{IEEEbiography}[{\includegraphics[width=1in,height=1.25in,clip,keepaspectratio]{Figures_v2/karlphoto.jpg}}]{Karl Raymond Dela Cruz Roque}
(Graduate Student Member, IEEE) is with the Graduate School of Engineering, Department of Electrical Engineering, in Nagoya University, Japan. 
He received the B.Eng. Degree in Automotive Engineering with a Specialization in Electrical Engineering from Nagoya University in Nagoya, Japan in 2024. He is currently pursuing the M.Eng. Degree in Electrical Engineering at Nagoya University, Nagoya, Japan. Since 2023, he has been a Research Assistant at the Nagoya University Yamamoto Laboratory for Power Electronics. His research interests include the design of magnetic components, particularly integrated magnetics and thermal management, as well as resonant converters, control design, and system modeling.
\end{IEEEbiography}
%
% Placeholder for other biographies if needed
\vfill

\end{document}